% USERS_MANUAL_JA.tex — xinu-rpi4 ユーザーズマニュアル (日本語版)
% Build: lualatex USERS_MANUAL_JA.tex   (目次のため 2 回実行)
\documentclass[a4paper,11pt]{ltjsarticle}

\usepackage{luatexja-fontspec}
\setmonofont{Menlo}[Scale=0.85]    % コード/ツリーの box-drawing と矢印用
\usepackage[a4paper,margin=2.0cm]{geometry}
\usepackage{amssymb}
\usepackage{xcolor}
\usepackage{fancyvrb}
\usepackage{longtable}
\usepackage{booktabs}
\usepackage{array}
\usepackage{enumitem}
\usepackage{newunicodechar}
\usepackage[hidelinks]{hyperref}

% --- 本文中で使う Unicode 記号 (コードブロックは Menlo が直接描画) ---
\newunicodechar{→}{\textrightarrow}
\newunicodechar{↔}{\ensuremath{\leftrightarrow}}
\newunicodechar{×}{\ensuremath{\times}}
\newunicodechar{…}{\dots}
\newunicodechar{≥}{\ensuremath{\geq}}
\newunicodechar{≠}{\ensuremath{\neq}}
\newunicodechar{✅}{\textcolor{green!55!black}{\checkmark}}
\newunicodechar{⏳}{\textcolor{orange!80!black}{\textbf{\small WIP}}}
\newunicodechar{❌}{\textcolor{red!70!black}{\ensuremath{\times}}}
\newunicodechar{⚠}{\textcolor{orange!80!black}{\textbf{!}}}

% --- コードブロック ---
\DefineVerbatimEnvironment{code}{Verbatim}%
  {fontsize=\small,frame=single,framerule=0.3pt,rulecolor=\color{gray!60},%
   framesep=5pt,xleftmargin=2pt,xrightmargin=2pt}

\setlength{\parindent}{0pt}
\setlength{\parskip}{4pt}
\renewcommand{\arraystretch}{1.25}

\title{\textbf{Xinu \texttt{xinu-rpi4} ユーザーズマニュアル}\\[2pt]
       \large Raspberry Pi 4 (BCM2711 / Cortex-A72, AArch64) 向け Embedded Xinu}
\author{}
\date{}

\begin{document}
\maketitle
\vspace{-2.2em}
\begin{center}\rule{0.9\textwidth}{0.4pt}\end{center}

Raspberry Pi 4 (BCM2711 / Cortex-A72, AArch64) 向けの Embedded Xinu 移植版。
同じソースツリーから QEMU \texttt{virt} 機にもクロスビルドできる。本マニュアルは
「動かす側」向け --- リポジトリの開発ロードマップは \texttt{README.md} を参照のこと。

{\small\tableofcontents}
\bigskip

%======================================================================
\section{対象ターゲットとビルド成果物}

\begin{center}\footnotesize
\begin{tabular}{@{}l l l l@{}}
\toprule
\textbf{ターゲット} & \textbf{ボード / 機種} & \textbf{出力イメージ} & \textbf{UART0 base}\\
\midrule
\verb|pi4|  & Raspberry Pi 4 (BCM2711) & \verb|kernel8.img|     & \verb|0xFE201000|\\
\verb|qemu| & QEMU \verb|-M virt -cpu cortex-a76| & \verb|kernel_virt.img| & \verb|0x09000000|\\
\bottomrule
\end{tabular}
\end{center}

load address は \verb|0x80000| (Pi 4) / \verb|0x40080000| (QEMU)。2 つのイメージは
同一ソースから \verb|-D<TARGET>| と異なるリンカスクリプトで作り分けている。

\subsection{機能対応マトリクス (重要)}

「全ターゲットで動く」訳ではない。\textbf{実機 (Pi 4) でしか動かない機能} が
複数あるので要注意:

\begin{center}\footnotesize
\begin{tabular}{@{}l c c p{6.4cm}@{}}
\toprule
\textbf{サブシステム} & \textbf{Pi 4} & \textbf{QEMU} & \textbf{備考}\\
\midrule
PL011 UART0 シェル        & ✅ & ✅ & 全機種で動く土台\\
協調的スケジューラ (S0)   & ✅ & ✅ & \verb|procdemo| / \verb|pingpong|\\
プリエンプティブ (S1)     & ✅ & ⏳ & Pi 4 のみ 100 Hz tick 稼働 (\verb|ticks|)\\
HDMI フレームバッファ     & ✅ & — & Pi 4 で実用\\
ウィンドウマネージャ      & ✅ & — & 1280×960 仮想デスクトップは Pi 4 のみ\\
SD カード + FAT32        & ✅ & — & EMMC ドライバが BCM2711 向け\\
Ethernet (GENET)         & ✅ & — & NET-E 完了 (BCM2711 GENET)\\
USB (キーボード/HID)      & ✅ & — & VL805 PCIe xHCI 経由 (USB-2.0 ポート)\\
DHCP / TCP               & ❌ & — & コードは在るが dispatch OFF (§7.4)\\
\bottomrule
\end{tabular}
\end{center}

⚠~ソースツリーには \texttt{device/genet/}, \texttt{device/sd/},
\texttt{device/usb/xhci/} 等が同居しており、これらは \textbf{Pi 4 (BCM2711) の
レジスタアドレスで初期化} される。\verb|-DGENET_BASE| 等は Pi 4 ビルドでのみ定義
され、QEMU ビルドでは undef になって対応コードがリンクから外れる作り。

\subsection{起動シーケンス (特殊な部分)}

\texttt{loader/boot.S} は通常の bare-metal stub よりひと工夫が要る:

\begin{enumerate}[leftmargin=1.6em,itemsep=2pt]
\item \textbf{Linux ARM64 Image ヘッダが必須。} Pi 4 の EEPROM bootloader は
  ファイル先頭から \verb|0x38| オフセットに \verb|"ARM\x64"| マジックが無いと
  カーネルへジャンプしないケースがある。\texttt{boot.S} の \verb|_start:| 直下に
  64 バイトのヘッダを配置。無いと \textbf{HDMI に虹色が出たまま固まる}。
\item \textbf{MPIDR\_EL1 によるコア固定。} boot core 以外 (core1/2/3) を
  \verb|wfe| ループに park。
\item \textbf{スタック設定。} \verb|sp = _start| (= \verb|0x80000|)。
\item \textbf{DTB アドレス退避。} ファームウェアから \verb|x0| に渡される DTB の
  物理アドレスを BSS clear の間退避し \verb|.data dtb_addr| に格納。
\item \textbf{BSS クリア → \texttt{kernel\_main}。} \verb|kernel_main| から戻れば
  secondary と同じ \verb|wfe| パークへ落ちる。
\end{enumerate}

EL のドロップ (EL2\,\textrightarrow\,EL1) は自分で実装していない (QEMU も同じ
前提)。\verb|armstub| や \verb|kernel_old=1| を入れると状況が変わる。

%======================================================================
\section{用意するもの}

\subsection{ハードウェア (実機の場合)}
\begin{itemize}[leftmargin=1.4em,itemsep=1pt]
\item Raspberry Pi 4 (BCM2711)
\item microSD カード (FAT32 で \verb|bootfs| パーティション)。一番簡単なのは
  Raspberry Pi OS を焼き、後で \verb|kernel8.img| と \verb|config.txt| を上書き
\item USB-シリアル変換ケーブル (3.3V、115200 8N1) --- ヘッダピン 8
  (TXD\,\textrightarrow\,GPIO14) / 10 (RXD\,\textrightarrow\,GPIO15) / 6 (GND)
\item ネットワーク機能を使う場合: イーサネットケーブル
\end{itemize}

\subsection{ホスト側ソフトウェア}
macOS / Linux と AArch64 クロスツールチェイン (どちらか):
\begin{code}
brew install aarch64-elf-gcc           # Homebrew (推奨)
brew install --cask gcc-arm-embedded   # ARM 公式
\end{code}
QEMU で動かすなら \verb|qemu-system-aarch64| (\verb|brew install qemu|)、および
ターミナルシリアルクライアント (\verb|screen|、\verb|minicom|、\verb|picocom|)。

%======================================================================
\section{ビルド手順}

ビルドは必ず \verb|compile/| ディレクトリで行うこと (\verb|Makefile| が
\verb|VPATH| で \verb|../loader/|、\verb|../device/...| を集める作り)。

\subsection{基本コマンド}
\begin{code}
cd /Users/kodamay/projects/xinu-rpi4/compile

make pi4            # -> compile/kernel8.img      (Raspberry Pi 4 実機)
make qemu           # -> compile/kernel_virt.img  (QEMU virt)
make                # = make all = pi4 + qemu
make clean          # オブジェクト・全 .img を削除
\end{code}
各ターゲットのオブジェクトは別ツリー (\verb|obj/pi4/|, \verb|obj/qemu/|) に
分離され、続けて別ターゲットをビルドしても再コンパイルは差分のみ。

\subsection{ツールチェインの自動検出}
\verb|$(GCCPATH)/bin/aarch64-elf-gcc| → \verb|aarch64-none-elf-gcc| の順に探す。
既定 \verb|GCCPATH| は \verb|/opt/homebrew|。別の場所なら:
\begin{code}
make pi4 GCCPATH=$HOME/aarch64/arm-gnu-toolchain-...-aarch64-none-elf
make pi4 CROSS=aarch64-linux-gnu-
\end{code}

\subsection{主な Make 変数}
\begin{center}\footnotesize
\begin{tabular}{@{}l l p{7.0cm}@{}}
\toprule
\textbf{変数} & \textbf{既定値} & \textbf{用途}\\
\midrule
\verb|GCCPATH| & \verb|/opt/homebrew| & クロスコンパイラのインストール先\\
\verb|CROSS|   & (自動検出)          & プレフィックス (\verb|aarch64-elf-| 等)\\
\verb|SDCARD|  & \verb|/Volumes|     & \verb|make install_*| のマウント親\\
\verb|DEST|    & \verb|$(SDCARD)/bootfs| & コピー先パーティション\\
\bottomrule
\end{tabular}
\end{center}

\subsection{ビルド時の主な定義 (CFLAGS で渡る)}
\begin{center}\footnotesize
\begin{tabular}{@{}l l l@{}}
\toprule
\textbf{マクロ} & \textbf{Pi 4} & \textbf{QEMU virt}\\
\midrule
\verb|-mcpu|       & \verb|cortex-a72| & \verb|cortex-a76|\\
\verb|UART0_BASE|  & \verb|0xFE201000| & \verb|0x09000000|\\
\verb|MBOX_BASE|   & \verb|0xFE00B880| & (未定義)\\
\verb|SD_BASE|     & \verb|0xFE340000| & ---\\
\verb|USB_BASE|    & \verb|0xFE980000| & ---\\
\verb|GIC_BASE|    & \verb|0xFF840000| & ---\\
\verb|PCIE_BASE|   & \verb|0xFD500000| & ---\\
\verb|GENET_BASE|  & \verb|0xFD580000| & ---\\
\verb|HEAP_END|    & \verb|0x40000000| (1G) & \verb|0x50000000|\\
\verb|SKIP_MBOX|   & ---               & 定義 (簡略化)\\
\verb|KERNEL_NAME| & \verb|"kernel8.img"| & \verb|"kernel_virt.img"|\\
\verb|BOARD_NAME|  & \verb|"Pi4"|      & \verb|"virt"|\\
\verb|SOC_NAME|    & \verb|"BCM2711"|  & \verb|"QEMU"|\\
\bottomrule
\end{tabular}
\end{center}

共通 CFLAGS:
\verb|-Wall -O2 -ffreestanding -nostdinc -nostdlib -nostartfiles -mgeneral-regs-only|。
大半のコードで FPU/NEON は使わず、\verb|cc.c| / \verb|llm.c| のみ FP 有効でビルドする。

\subsection{リンカスクリプト}
\begin{center}\footnotesize
\begin{tabular}{@{}l l l@{}}
\toprule
\textbf{スクリプト} & \textbf{対象} & \textbf{エントリ・ロードアドレス}\\
\midrule
\verb|compile/link.ld|      & Pi 4      & \verb|0x80000|\\
\verb|compile/link_virt.ld| & QEMU virt & \verb|0x40080000|\\
\bottomrule
\end{tabular}
\end{center}
どちらも \verb|.text.boot| を先頭に置き、\verb|__bss_start| / \verb|__bss_size|
(8-byte 単位) をエクスポートする leex 風の構成。

%======================================================================
\section{SD カードへの書き込み}

\subsection{Pi 4}
\begin{code}
cd compile
make install_pi4 SDCARD=/Volumes
# -> /Volumes/bootfs に kernel8.img と sdcard/config_pi4.txt をコピー
\end{code}
\verb|make install| は \verb|install_pi4| のエイリアス。

\subsection{手動で焼く場合}
\begin{code}
diskutil mount /dev/disk4s1                # -> /Volumes/bootfs
cp compile/kernel8.img /Volumes/bootfs/
cp sdcard/config_pi4.txt /Volumes/bootfs/config.txt
sync
diskutil eject /Volumes/bootfs
\end{code}
⚠~\verb|kernel8.img| の md5 が SD カード上のものと一致するか必ず確認すること。
古いイメージで起動して混乱するのは「あるある」。

\subsection{SD カードに必要なもの}
\verb|kernel8.img| と \verb|config.txt| だけでは起動しない。同じ FAT32 パーティ
ションに Raspberry Pi OS の \verb|bootcode.bin|, \verb|start4.elf|,
\verb|fixup4.dat| ほかファームウェア blob 一式も必要。Pi OS を焼いてから
\verb|kernel8.img| と \verb|config.txt| だけ上書きするのが最速。

%======================================================================
\section{シリアルコンソールへの接続}
\begin{code}
screen /dev/tty.usbserial-XXXX 115200                       # Mac
minicom -b 115200 -o -D /dev/tty.usbserial-XXXX             # Mac
sudo screen /dev/ttyUSB0 115200                             # Linux
\end{code}
\verb|screen| を終了するときは \verb|Ctrl-a k| → \verb|y|。電源投入後、約 5 秒
以内に次のバナーが出る:
\begin{code}
================================================
  Xinu Pi4 hello (AArch64, BCM2711, kernel8.img)
  PL011 UART0 @ 0xFE201000, 115200 8N1
  bootstrap: leex-style stub + xinu-rpi4 main
================================================

Round 1 phase B/U done -- entering interactive shell.
type `help` for the command list.
xinu-pi4$ _
\end{code}

%======================================================================
\section{シェルコマンドリファレンス}
\verb|help| または \verb|?| で常に最新の一覧が見られる。現在のコマンド一覧:
\begin{center}\footnotesize
\begin{longtable}{@{}l p{8.6cm}@{}}
\toprule
\textbf{コマンド} & \textbf{用途}\\
\midrule \endhead
\verb|help| / \verb|?|    & 登録コマンド一覧を表示\\
\verb|echo <words>|       & 引数をそのままエコーバック\\
\verb|hello|             & 動作確認用の挨拶\\
\verb|mem|               & \verb|__bss_start| / \verb|__bss_end| / \verb|_end| を表示\\
\verb|peek <hex_addr>|    & 32-bit MMIO ワードを読む (例: \verb|peek 0xfe201018|)\\
\verb|uptime|           & 生の \verb|CNTPCT_EL0| (汎用タイマカウンタ)\\
\verb|ticks|            & 100 Hz タイマ tick の累積数 (S1)\\
\verb|ps|               & コア / EL ステータス簡易表示\\
\verb|halt|             & DAIF マスク + PSCI \verb|SYSTEM_OFF| (QEMU はクリーン終了)\\
\verb|reboot|           & スタブ — 電源再投入待ちでスピン\\
\verb|pwd ls cd mkdir rmdir touch cat write edit rm cp mv tree| & RAM 上のファイルシステム (揮発性。§6.3)\\
\verb|cc <file.c>|       & C をオンデバイスでコンパイル\&実行 (JIT \textrightarrow{} AArch64。§6.3)\\
\verb|aload| / \verb|amsg| & 常駐 AIPL アクタのロード / メッセージ送信 (§6.3)\\
\verb|actordemo| / \verb|selectdemo| & アクタ ping-pong / 名前付き選択受信 (§6.3)\\
\verb|vmtest| / \verb|vmdemand| & VA≠PA リマップ / デマンドページング仮想メモリ (§6.4)\\
\verb|llm [prompt]|      & 組込み LLM でテキスト生成 (§6.3)\\
\verb|preempt|          & タイマ駆動プリエンプティブ RR のデモ\\
\verb|pingpong [N]|      & AIPL 風 2 アクター協調 PingPong (N=1..50)\\
\verb|procdemo [N]|      & 実 ctxsw による 2 プロセスデモ (N=1..30)\\
\verb|usb|              & xHCI / DWC2 USB 診断 (Pi 4 のみ)\\
\verb|wifi ...|          & WiFi + メッシュ: \verb|probe|/\verb|scan|/\verb|up|/\verb|adhoc|/\verb|aodv|/\dots\ (§7.5--§7.6)\\
\verb|rxstat| / \verb|tcpstat| & RX リング / TCP listener カウンタ\\
\verb|pan <dx> <dy>| \verb|view| \verb!autopan [on|off]! & ウィンドウマネージャのビューポート操作\\
\bottomrule
\end{longtable}
\end{center}

\subsection{procdemo の見どころ}
\begin{code}
xinu-pi4$ procdemo 3
procdemo: created pid=1 (ping) and pid=2 (pong), iters=3
  [Ping pid=1] tick 1
  [Pong pid=2] tock 1
  [Ping pid=1] tick 2
  [Pong pid=2] tock 2
  [Ping pid=1] tick 3
  [Pong pid=2] tock 3
  [Ping pid=1] exit at iter 3
  [Pong pid=2] exit at iter 3
procdemo: both processes exited; back in shell.
\end{code}
\verb|pid=N| はグローバル \verb|currpid| から実行時に読み出した値なので、列方向の
交互パターン自体がスケジューラが本当に切替えている証拠になる。スタックは
\verb|proc_exit| では回収していない (S1 で IRQ tick 経由で回収予定)。

\subsection{pingpong との違い}
\begin{center}\footnotesize
\begin{tabular}{@{}l p{4.2cm} p{5.2cm}@{}}
\toprule
\textbf{観点} & \verb|pingpong| & \verb|procdemo|\\
\midrule
アクター   & 静的 2 つ (\verb|Ping|, \verb|Pong|) & \verb|proc_create| した実プロセス\\
スイッチ   & 単一スタック上のディスパッチ & \verb|ctxsw.S| による実 callee-saved 保存\\
停止条件   & 両者の inbox が空 & 両者が \verb|proc_exit()|\\
\bottomrule
\end{tabular}
\end{center}

\subsection{オンデバイス・ツール (ファイルシステム / C JIT / アクタ / LLM)}
上記のデモに加え、シェルには自己完結したツールチェインが載っており、ボード上で
コードを書いて・コンパイルして・実行できます。
\begin{itemize}[leftmargin=1.4em,itemsep=1pt]
\item \textbf{RAM 上のファイルシステム} --- \verb|pwd| \verb|ls| \verb|cd|
  \verb|mkdir| \verb|rmdir| \verb|touch| \verb|cat| \verb|write| \verb|edit|
  \verb|rm| \verb|cp| \verb|mv| \verb|tree| がメモリ上の小さなツリーを操作
  (揮発性。再起動で消える)。
\item \textbf{C JIT (\texttt{cc})} --- \verb|cc <file.c>| が C のサブセットを
  ネイティブ AArch64 にコンパイルしてその場で実行。同じコンパイラは HTTP の
  \verb|POST /compile| (本文 = C ソース) からも使えます。
\item \textbf{AIPL アクタ} --- \verb|aload <file.c>| で常駐アクタをロード、
  \verb|amsg <actor> <method> [arg]| で送信。\verb|actordemo| は 2 アクターの
  ping-pong を実 Xinu プロセスとして実行、\verb|selectdemo| はガード付き受信を
  示します。HTTP の \verb|/actor|・\verb|/send|・\verb|/gc| がアクタ一覧・送信・
  アクタプール GC を公開。
\item \textbf{組込み LLM (\texttt{llm})} --- 小さな transformer がイメージに焼き
  込まれており、\verb|llm [prompt]| でオンデバイス生成 (HTTP の \verb|/chat| でも)。
\end{itemize}

\subsection{仮想メモリ (\texttt{vmtest} / \texttt{vmdemand})}
カーネルは MMU 有効 (identity map) + VA \verb|0x80000000|..\verb|0x80400000|
(4 MiB, 512 フレームプール) の\textbf{デマンドページング窓}で動作します:
\begin{itemize}[leftmargin=1.4em,itemsep=1pt]
\item \verb|vmtest| --- 同じ物理ページを別の仮想アドレスにマップして変換を実証
  (VA ≠ PA)。
\item \verb|vmdemand| --- デマンド窓の 64 ページに触れる。初回は \verb|0x40| 回の
  ページフォルト (1 ページ 1 回) でリードバック OK、2 回目は \verb|0x0| フォルト
  (既にマップ済み)。フォルトカウンタは \verb|/fault| で確認。
\end{itemize}
ローカルでもリモートでも実行可: 例
\verb|curl "http://192.168.3.100/shell?cmd=vmdemand"|。

%======================================================================
\section{ネットワーク機能 (Pi 4)}
ネットワークドライバは Pi 4 GENET (BCM2711) で動作する。

\textbf{動作する機能:} ARP request/reply、ICMP echo (ping)、静的 IP/MAC
\verb|192.168.3.100| / \verb|d8:3a:dd:a7:fd:bf| (\verb|loader/main.c:707| 付近で
設定)、生 broadcast/unicast 送信、RX リング 16-slot。

\subsection{Mac から ping する手順}
\begin{code}
sudo arp -s 192.168.3.100 d8:3a:dd:a7:fd:bf
ping 192.168.3.100
\end{code}
RTT は約 900\,ms (ウィンドウマネージャの tick 律速)。シェルからは \verb|rxstat| /
\verb|ticks| で状況確認。

\subsection{DHCP / TCP の現状}
\verb|system/dhcp_client.c| と \verb|system/tcp_server.c| はツリーに含まれるが、
\verb|rx_tick| からの dispatch を \textbf{OFF} にしている: 連続 broadcast TX で
GENET TX ring が劣化し ICMP が止まる、また \verb|tcp_handle_packet()| を
\verb|rx_tick| から呼ぶと以降の ICMP echo reply が止まるため。再有効化は今後の
スプリント (NET-G bisect) で対応予定。

\subsection{WiFi (BCM43455)}
Pi 4 はオンボードの BCM43455 WiFi チップを持ちます (上記の有線 GENET とは別系統)。
操作はすべてシリアルシェル (\verb|xinu-pi4$|) から行います: ファームウェアを起こし、
スキャンし、接続します。\textbf{起動時に自動接続はしません。}
\begin{center}\footnotesize
\begin{tabular}{@{}l p{8.2cm}@{}}
\toprule \textbf{コマンド} & \textbf{説明}\\ \midrule
\verb|wifi probe| & M0/M1 bring-up: チップにファームウェアを DL (起動ごとに 1 回)\\
\verb|wifi scan| & 近隣 AP を escan\\
\verb|wifi up <ssid> <pass>| & join (WPA2-PSK) + DHCP。常駐 ARP/ICMP 応答を開始\\
\verb|wifi join <ssid> <pass>| & join のみ (DHCP なし)\\
\verb|wifi dhcp| & DHCP リース取得\\
\verb|wifi off| & 無線停止 / 切断\\
\verb|wifi ping <a.b.c.d> [count]| & ICMP echo クライアント\\
\verb|wifi serve [secs]| & ARP/ICMP 応答を N 秒間実行 (既定 30)\\
\verb|wifi resolve <host>| / \verb|wifi web <host>| & DNS / DNS + HTTP GET\\
\verb|wifi ntp [a.b.c.d]| & NTP 時刻クライアント\\
\verb|wifi tftp <ip> <file>| / \verb|wifi netboot <ip> <file>| & TFTP 取得 / 取得 + チェインロード\\
\bottomrule
\end{tabular}
\end{center}
\begin{code}
xinu-pi4$ wifi probe          # ファームウェア DL (起動ごとに 1 回)
xinu-pi4$ wifi scan           # 周囲の AP を確認
xinu-pi4$ wifi up MyHome-5G mypassword
wifi up: connected; ARP/ICMP responder is now persistent.
xinu-pi4$ wifi ping 8.8.8.8
\end{code}
再起動後は自動再接続しません。起動のたびに \verb|wifi probe| + \verb|wifi up| を
実行してください。\verb|wifi netboot <ip> <file>| は WiFi 越しにカーネルを取得して
チェインロードします (有線 chainload と同様、SD 交換なしの更新経路)。

\subsection{複数 Xinu でのメッシュネットワーク (MANET ad-hoc)}
複数の Pi 4 (および Pi 3) Xinu ボードを \textbf{アクセスポイント無し} で直接つなぎ、
ピアツーピアのメッシュを組めます。802.11 の IBSS (ad-hoc) モード + オンデマンドの
AODV ルーティングを使います (ドローン HIL デモの Pi 4 + Pi 3 ノードと同じ MANET
スタック)。各ノードは同じ ad-hoc セルに、別々のノード番号で参加します:
\begin{code}
xinu-pi4$ wifi probe                       # 起動ごとに 1 回
xinu-pi4$ wifi adhoc <cell-ssid> [ch] [node]
\end{code}
\begin{itemize}[leftmargin=1.4em,itemsep=1pt]
\item \verb|<cell-ssid>| --- ad-hoc セル名。\textbf{全ノードで同じ名前・同じ
  チャンネル}にすること。
\item \verb|[ch]| --- 802.11 チャンネル (既定 6)。
\item \verb|[node]| --- 自分のノード番号 (既定 1)。静的 IP \verb|10.0.0.<node>/24|
  を得ます。\textbf{ノードごとに別の番号}を与えること。
\end{itemize}
例 --- チャンネル 6 の 3 ノードセル:
\begin{code}
xinu-pi4$ wifi adhoc mesh1 6 1      # -> 10.0.0.1
xinu-pi4$ wifi adhoc mesh1 6 2      # -> 10.0.0.2
xinu-pi4$ wifi adhoc mesh1 6 3      # -> 10.0.0.3
\end{code}
電波の届く範囲のノードは \verb|10.0.0.x| で直接通信できます (\verb|wifi ping 10.0.0.2|)。

\textbf{マルチホップ (AODV)。} 宛先が直接届かないとき、\verb|wifi aodv <ip>| で
オンデマンドに経路探索します:
\begin{code}
xinu-pi4$ wifi aodv 10.0.0.3
[wifi] === AODV discover 10.0.0.3 ===
[wifi] aodv: RREQ id=1 for 10.0.0.3
[wifi] *** AODV route: 10.0.0.3 via 10.0.0.2, 2 hop ***
\end{code}
最小 AODV モジュール (M13, RFC 3561 コア) が RREQ をブロードキャスト (UDP port 654)
し、範囲内の各ノードが中継して逆経路を記録、宛先が RREP を返して順経路を確立します。
各ノードはピアのために自動で中継するので、直接届かない宛先も中間ノードが転送します
(経路表は最大 16 エントリ)。IBSS/ad-hoc はインフラ (\verb|wifi up|) モードとは独立
です (AP に繋がっている場合は先に \verb|wifi off|)。再起動後は復元しないので、各
ノードで \verb|wifi probe| + \verb|wifi adhoc| を再実行してください。

\subsection{HTTP 制御面とリモートシェル}
\verb|system/tcp_server.c| は有線インタフェース上で HTTP サーバ (既定ポート 80、
\verb|192.168.3.100|) を動かし、ボードをネットワーク越しに丸ごと操作できます:
\begin{center}\footnotesize
\begin{tabular}{@{}l p{8.0cm}@{}}
\toprule \textbf{ルート} & \textbf{内容}\\ \midrule
\verb|GET /shell?cmd=<cmd>| & 任意のシェルコマンドを実行し出力を返す (stdin なし)\\
\verb|POST /compile| & C を JIT コンパイル\&実行 (本文 = ソース)\\
\verb|GET /chat| & オンデバイス LLM チャット\\
\verb|GET /actor| \verb|/send| \verb|/gc| \verb|/jitstats| & アクタ一覧 / 送信 / GC / JIT カウンタ\\
\verb|GET /usbdiag| \verb|/pcie-init| \verb|/pcie-enum| \verb|/xhci-reset| & USB / PCIe bring-up 診断\\
\verb|GET /fault| \verb|/mmio-read| \verb|/mmio-write| \verb|/mmio-sweep| & フォルトカウンタ + 生 MMIO peek/poke\\
\verb|POST /reboot| & BCM2711 watchdog リセット\\
\verb|POST /chainload| & 新カーネルをアップロードして起動 (SD 交換なし)\\
\bottomrule
\end{tabular}
\end{center}
実行時診断にこれら HTTP カウンタを使うのは、\verb|uart_puts| が HTTP ワーカ
コンテキストでデッドロックするためです (bring-up ログはブート時にシリアルへ)。

\textbf{SD 交換なしのカーネル更新 (chainload)。} \verb|POST /chainload| は新カーネル
をアップロードして RAM 上でジャンプします。アップロードはヘルパスクリプトが包みます:
\begin{code}
python3 tools/remote_chainload.py 192.168.3.100 compile/kernel8.img
\end{code}
不良イメージでも電源を入れ直すだけで戻ります (SD は無傷) ので開発ループが速く
なります。アップロードには HTTP サーバが応答可能であることが必要です。

%======================================================================
\section{QEMU で試す}
\begin{code}
cd compile
make qemu          # インタラクティブ起動。Ctrl-A X で抜ける
make qemu-smoke    # 既定のスクリプトを流し込み qemu-smoke.log に保存
\end{code}
QEMU 上では: \verb|MIDR_EL1| が \verb|0x414fd0b1| (QEMU が公開する Cortex-A76 の
part 番号。実機 Pi 4 では Cortex-A72)、\verb|halt| は PSCI SYSTEM\_OFF が
\verb|virt| でハンドルされクリーン終了、ネットワーク・USB・SD・HDMI 出力は無し。

%======================================================================
\section{config.txt の主な設定}
\verb|sdcard/config_pi4.txt| (\verb|make install_pi4| がカード上の
\verb|config.txt| へコピー) の重要項目:
\begin{center}\footnotesize
\begin{tabular}{@{}l p{7.2cm}@{}}
\toprule
\textbf{行} & \textbf{意味}\\
\midrule
\verb|arm_64bit=1|              & AArch64 で起動\\
\verb|kernel=kernel8.img|       & Pi 4 用カーネル名\\
\verb|kernel_address=0x80000|   & カーネルロードアドレス\\
\verb|enable_uart=1|            & UART を有効化\\
\verb|uart_2ndstage=1|          & bootloader 第 2 段の UART ログを出す\\
\verb|dtparam=uart0=on|         & UART0 を GPIO14/15 にルーティング\\
\verb|init_uart_clock=48000000| & PL011 基準クロックを 48 MHz に固定 (115200 用)\\
\verb|hdmi_force_hotplug=1|     & HDMI ホットプラグを強制 ON\\
\verb|framebuffer_width=640|    & ファームウェア確保 FB の幅\\
\verb|framebuffer_height=480|   & FB の高さ\\
\verb|framebuffer_depth=32|     & FB の bpp\\
\bottomrule
\end{tabular}
\end{center}
KMS overlay (\verb|dtoverlay=vc4-kms-v3d|) を残してあるが、Xinu 側はファームウェアが
用意した simple-framebuffer を mailbox 経由で借りる構成。

%======================================================================
\section{USB / 入力デバイスのサポート状況}
Pi 4 のフルサイズ USB-A ポートは \textbf{VL805 (PCIe xHCI)} 経由でぶら下がって
おり、xHCI ドライバを自作して USB キーボード・マウスを動かしている。
\begin{itemize}[leftmargin=1.4em,itemsep=1pt]
\item \textbf{USB-2.0 (黒) ポートに挿すこと。} USB-2.0 ハブ (Genesys、TT 動作)
  経由なので LS の割り込み IN split transaction が届く。USB-3.0 (青) は VIA USB3
  ハブ側に入り、周期転送が届かず動かない。
\item \verb|usb| コマンドは DWC2 (USB-C OTG) HCD の診断ダンプ専用。
\item VL805 のファームウェアは SD ブート時 PCIe RC がリセットのままになるため、
  \verb|NOTIFY_XHCI_RESET| mailbox tag (\verb|0x00030058|, devid \verb|0x100000|)
  を boot 時に発行して xHCI を起こしている (詳細は \verb|NEXT_SESSION_PI4.md|)。
\end{itemize}

\begin{center}\footnotesize
\begin{tabular}{@{}l l l@{}}
\toprule
\textbf{機種} & \textbf{入力経路} & \textbf{出力経路}\\
\midrule
Pi 4  & USB-A キーボード / UART0 シリアル & UART0 シリアル + HDMI FB\\
QEMU  & qemu のシリアル stdio          & qemu のシリアル stdio\\
\bottomrule
\end{tabular}
\end{center}

%======================================================================
\section{トラブルシューティング}
\begin{center}\footnotesize
\begin{tabular}{@{}p{5.0cm} p{8.4cm}@{}}
\toprule
\textbf{症状} & \textbf{確認ポイント}\\
\midrule
シリアルに何も出ない            & \verb|screen| のデバイス名と速度、TX/RX 入れ違い、GND\\
バナーは出るがすぐ止まる        & \verb|kernel8.img| の md5 が古い。再ビルド + 再コピー\\
キー入力が echo されない        & ターミナルの local echo は OFF が正しい\\
\verb|peek| で 0xFFFFFFFF        & そのアドレス帯が未マップ。MMU は flat ID\\
USB キーボード/マウスが効かない  & USB-2.0 (黒) ポートに挿しているか (§10)\\
ping が通らない (Pi 4)          & \verb|sudo arp -s ...| 済か、ケーブル、\verb|rxstat| 増分\\
\verb|make qemu| が失敗          & \verb|qemu-system-aarch64| のバージョン ($\geq$11)\\
\verb|make install| permission   & 内部で \verb|sudo cp| するためパスワード入力\\
起動後すぐ reboot ループ        & \verb|config.txt| の \verb|kernel=| 行と実ファイル名\\
HDMI に虹色が残る               & ARM64 Image ヘッダ欠落の可能性。再ビルド\\
\bottomrule
\end{tabular}
\end{center}
ハマったときの定番動作:
\begin{code}
cd compile && make clean && make pi4
make install_pi4 SDCARD=/Volumes
diskutil eject /Volumes/bootfs
# SD カード抜き挿し → Pi の電源再投入
\end{code}

%======================================================================
\section{ディレクトリ構成 (利用者目線)}
\begin{code}
xinu-rpi4/
├── compile/                # ビルドはここで `make`
│   ├── Makefile
│   ├── kernel8.img         # Pi 4 用 (make pi4 で生成)
│   └── kernel_virt.img     # QEMU 用 (make qemu で生成)
├── sdcard/
│   └── config_pi4.txt      # Pi 4 用 (install_pi4 がコピー)
├── README.md               # 開発者向けロードマップ + 内部実装
├── USERS_MANUAL_JA.md      # 本マニュアルの Markdown 版
├── USERS_MANUAL_EN.md      # 英語版
└── NEXT_SESSION_PI4.md     # 進行中作業のハンドオフメモ
\end{code}

%======================================================================
\section{よくある操作レシピ集}
\textbf{ビルドから Pi 4 で起動するまで一気通貫:}
\begin{code}
cd /Users/kodamay/projects/xinu-rpi4/compile
make clean && make pi4
make install_pi4 SDCARD=/Volumes
diskutil eject /Volumes/bootfs
screen /dev/tty.usbserial-XXXX 115200
\end{code}
\textbf{コマンドを手早く確かめる (QEMU):}
\begin{code}
cd compile && make qemu
xinu-pi4$ procdemo 5
xinu-pi4$ pingpong 3
xinu-pi4$ halt          # Ctrl-A X で QEMU から脱出
\end{code}

%======================================================================
\section{参考リンク}
\begin{itemize}[leftmargin=1.4em,itemsep=1pt]
\item 上流 Xinu (32-bit): \url{https://github.com/yaskodama/xinu-rpi}
\item AArch64 boot stub の元: \url{https://github.com/radlyeel/leex}
\item シェル centry パターンの元: \url{https://github.com/davidxyz/xinuPi}
\item 開発ロードマップ (Round 1 計画書): abclcp-project リポジトリの
  \verb|aice-pi-evolution/experiments/| 配下
\end{itemize}

\section{ライセンス}
上流 Xinu と leex の BSD 系ライセンスを継承。詳細は \verb|LICENSE| を参照 (整備中)。

\end{document}
